\ifdefined\gkver\else\def\gkver{student}\fi
\documentclass[\gkver]{gaokaozhenti}
\begin{document}

\chapter{2006年全国理综Ⅰ}
%% paperType: 理综物理部分

\begin{enumerate}
\item
%% number: 14
%% paperName: 2006年全国理综Ⅰ第 14 题 6 分
%% typeId: 1
%% type: 单选题
%% chapter: 原子和原子核
%% point: 人工核反应
%% method: 无
%% score: 6
%% degree: 850
%% duplicateId: 0
%% body:
某原子核 \ce{^{A}_{Z}X} 吸收一个中子后，放出一个电子，分裂为两个 $\alpha$ 粒子。由此可知\xzanswer{A}

\fourchoices[answer=A]
{$A=7$，$Z=3$}
{$A=7$，$Z=4$}
{$A=8$，$Z=3$}
{$A=8$，$Z=4$}

%% answer: A
%% memo:
\memoanswer{
\ce{^{A}_{Z}X} 吸收中子后变为 \ce{^{A+1}_{Z}X}，\ce{^{A+1}_{Z}X} 放出电子后变为 \ce{^{A+1}_{Z+1}X}，\ce{^{A+1}_{Z+1}X} 可分裂为两个 $\alpha$ 粒子，则 $A+1=8$，$Z+1=4$，所以 $A=7$，$Z=3$，故 A 正确。
}

\item
%% number: 15
%% paperName: 2006年全国理综Ⅰ第 15 题 6 分
%% typeId: 1
%% type: 单选题
%% chapter: 光学
%% point: 光子说
%% method: 无
%% score: 6
%% degree: 850
%% duplicateId: 0
%% body:
红光和紫光相比\xzanswer{B}

\fourchoices[answer=B]
{红光光子的能量较大；在同一种介质中传播时红光的速度较大}
{红光光子的能量较小；在同一种介质中传播时红光的速度较大}
{红光光子的能量较大；在同一种介质中传播时红光的速度较小}
{红光光子的能量较小；在同一种介质中传播时红光的速度较小}

%% answer: B
%% memo:
\memoanswer{
红光的折射率比紫光的折射率小，波长比紫光波长要长，因此红光的频率较小，光子能量较小，又 $v=\frac{c}{n}$，所以在同一介质中红光的传播较快，故 B 正确。
}

\item
%% number: 16
%% paperName: 2006年全国理综Ⅰ第 16 题 6 分
%% typeId: 1
%% type: 单选题
%% chapter: 曲线运动
%% point: 人造地球卫星
%% method: 无
%% score: 6
%% degree: 650
%% duplicateId: 0
%% body:
我国将要发射一颗绕月运行的探月卫星“嫦娥 1 号”。设该卫星的轨道是圆形的，且贴近月球表面。已知月球的质量约为地球质量的 $\frac{1}{81}$，月球的半径约为地球半径的 $\frac{1}{4}$，地球上的第一宇宙速度约为 $7.9\Ukms$，则该探月卫星绕月运行的速率约为\xzanswer{B}

\fourchoices[answer=B]
{$0.4\Ukms$}
{$1.8\Ukms$}
{$11\Ukms$}
{$36\Ukms$}

%% answer: B
%% memo:
\memoanswer{
由 $G\frac{Mm}{r^{2}}=m\frac{v^{2}}{r}$ 得 $v=\sqrt{\frac{GM}{r}}$，又月球的质量约为地球质量的 $\frac{1}{81}$，月球的半径约为地球半径的 $\frac{1}{4}$，所以探月卫星绕月运行的速率约为地球上第一宇宙速度的 $\frac{2}{9}$，$v=7.9\Ukms\times\frac{2}{9}=1.76\Ukms$，故 B 正确。
}

\item
%% number: 17
%% paperName: 2006年全国理综Ⅰ第 17 题 6 分
%% typeId: 2
%% type: 多选题
%% chapter: 磁场
%% point: 洛伦兹力
%% method: 无
%% score: 6
%% degree: 650
%% duplicateId: 0
%% body:
图中为一“滤速器”装置示意图。a、b 为水平放置的平行金属板，一束具有各种不同速率的电子沿水平方向经小孔 O 进入 a、b 两板之间。为了选取具有某种特定速率的电子，可在 a、b 间加上电压，并沿垂直于纸面的方向加一匀强磁场，使所选电子仍能够沿水平直线 OO$'$ 运动，由 O$'$ 射出。不计重力作用。可能达到上述目的的办法是\xzanswer{AD}

\onepicture[w=6.0cm]{figs/17.png}

\fourchoices[answer=AD]
{使 a 板电势高于 b 板，磁场方向垂直纸面向里}
{使 a 板电势低于 b 板，磁场方向垂直纸面向里}
{使 a 板电势高于 b 板，磁场方向垂直纸面向外}
{使 a 板电势低于 b 板，磁场方向垂直纸面向外}

%% answer: AD
%% memo:
\memoanswer{
若 a 板电势高于 b 板，则电子受电场力向上，要使电子沿直线运动，则电子受到的洛伦兹力向下，由左手定则可知，磁场的方向垂直纸面向里，故 A 正确；若 a 板电势低于 b 板，则电子受电场力向下，要使电子沿直线运动，则电子受到的洛伦兹力向上，由左手定则可知，磁场的方向垂直纸面向外，故 D 正确。
}

\item
%% number: 18
%% paperName: 2006年全国理综Ⅰ第 18 题 6 分
%% typeId: 1
%% type: 单选题
%% chapter: 气体性质
%% point: 分子力
%% method: 无
%% score: 6
%% degree: 650
%% duplicateId: 0
%% body:
下列说法中正确的是\xzanswer{D}

\fourchoices[answer=D]
{气体的温度升高时，分子的热运动变得剧烈，分子的平均动能增大，撞击器壁时对器壁的作用力增大，从而气体的压强一定增大}
{气体体积变小时，单位体积的分子数增多，单位时间内打到器壁单位面积上的分子数增多，从而气体的压强一定增大}
{压缩一定量的气体，气体的内能一定增加}
{分子 a 从远外趋近固定不动的分子 b，当 a 到达受 b 的作用力为零处时，a 的动能一定最大}

%% answer: D
%% memo:
\memoanswer{
气体的压强从微观上来看与分子的平均动能和分子的密集程度有关，从宏观上来看与气体的温度和体积有关；气体升温时，分子的平均动能增大，但气体体积可能增大，所以气体压强不一定增大，故 A 错误；气体体积变小时，单位体积的分子数增多，但气体的温度可能降低，气体压强不一定增大，故 B 错误；压缩一定量的气体，外界对气体做功，但气体可能放热，气体内能不一定增加，故 C 错误；当分子 a 从远处趋近固定不动的分子 b 时，分子间开始表现为引力，分子力做正功，分子 a 的动能一直增加，当 a 到达受 b 的作用力为零处，即平衡位置时，此后分子间表现为斥力，分子力做负功，分子 a 的动能逐渐减小，所以当 a 到达受 b 的作用力为零处时，a 的动能一定最大，故 D 正确。
}

\item
%% number: 19
%% paperName: 2006年全国理综Ⅰ第 19 题 6 分
%% typeId: 2
%% type: 多选题
%% chapter: 机械振动
%% point: 受迫振动共振
%% method: 无
%% score: 6
%% degree: 650
%% duplicateId: 0
%% body:
一砝码和一轻弹簧构成弹簧振子，图 a 所示的装置可用于研究该弹簧振子的受迫振动。匀速转动把手时，曲杆给弹簧振子以驱动力，使振子做受迫振动。把手匀速转动的周期就是驱动力的周期，改变把手匀速转动的速度就可以改变驱动力的周期。若保持把手不动，给砝码一向下的初速度，砝码便做简谐运动，振动图线如图 b 所示。当把手以某一速度匀速转动，受迫振动达到稳定时，砝码的振动图线如图 c 所示。若用 $T_{0}$ 表示弹簧振子的固有周期，$T$ 表示驱动力的周期，$y$ 表示受迫振动达到稳定后砝码振动的振幅，则\xzanswer{AC}

\threepicture[hA=3.8cm, wB=4.4cm, wC=4.4cm,
	labelA=2006全国理综Ⅰ19a, labelB=2006全国理综Ⅰ19b, labelC=2006全国理综Ⅰ19c]
{figs/19a.png}
{figs/19b.png}
{figs/19c.png}

\fourchoices[answer=AC]
{由图线可知 $T_{0}=4\Us$}
{由图线可知 $T_{0}=8\Us$}
{当 $T$ 在 $4\Us$ 附近时，$y$ 显著增大；当 $T$ 比 $4\Us$ 小得多或大得多时，$y$ 很小}
{当 $T$ 在 $8\Us$ 附近时，$y$ 显著增大；当 $T$ 比 $8\Us$ 小得多或大得多时，$y$ 很小}

%% answer: AC
%% memo:
\memoanswer{
当把手不动时，弹簧振子的振动图象如图 b 所示，由图可知，弹簧振子的固有周期 $T_{0}=4\Us$，故 A 正确；当把手匀速转动时，振子做受迫振动，受迫振动的周期等于驱动力的周期，与振子的固有周期无关，当驱动力的周期接近振子的固有周期时，振子共振，振幅增大，当驱动力的周期比振子的固有周期小得多或大得多时，振幅都很小，由此可知 C 正确。
}

\item
%% number: 20
%% paperName: 2006年全国理综Ⅰ第 20 题 6 分
%% typeId: 1
%% type: 单选题
%% chapter: 运动和力
%% point: 动量和冲量
%% method: 无
%% score: 6
%% degree: 650
%% duplicateId: 0
%% body:
一位质量为 $m$ 的运动员从下蹲状态向上起跳，经 $\Delta t$ 时间，身体伸直并刚好离开地面，速度为 $v$。在此过程中\xzanswer{B}

\fourchoices[answer=B]
{地面对他的冲量为 $mv+mg\Delta t$，地面对他做的功为 $\frac{1}{2}mv^{2}$}
{地面对他的冲量为 $mv+mg\Delta t$，地面对他做的功为零}
{地面对他的冲量为 $mv$，地面对他做的功为 $\frac{1}{2}mv^{2}$}
{地面对他的冲量为 $mv-mg\Delta t$，地面对他做的功为零}

%% answer: B
%% memo:
\memoanswer{
运动员向上起跳的过程中，由动量定理知，合外力的冲量等于物体动量的变化，则 $I-mg\Delta t=mv$，所以地面对运动员的冲量为 $I=mv+mg\Delta t$，又运动员刚离开地面，在地面对运动员作用力的方向上没有位移，因此地面对运动员做的功为零，故 B 正确。
}

\item
%% number: 21
%% paperName: 2006年全国理综Ⅰ第 21 题 6 分
%% typeId: 1
%% type: 单选题
%% chapter: 电磁感应
%% point: 电磁感应的电量
%% method: 无
%% score: 6
%% degree: 450
%% duplicateId: 0
%% body:
如图，在匀强磁场中固定放置一根串接一电阻 $R$ 的直角形金属导轨 aOb（在纸面内），磁场方向垂直纸面朝里，另有两根金属导轨 c、d 分别平行于 Oa、Ob 放置。保持导轨之间接触良好，金属导轨的电阻不计。现经历以下四个过程：①以速率 $v$ 移动 $d$，使它与 Ob 的距离增大一倍；②再以速率 $v$ 移动 c，使它与 Oa 的距离减小一半；③然后，再以速率 $2v$ 移动 c，使它回到原处；④最后以速率 $2v$ 移动 $d$，使它也回到原处。设上述四个过程中通过电阻 $R$ 的电量的大小依次为 $Q_{1}$、$Q_{2}$、$Q_{3}$ 和 $Q_{4}$，则\xzanswer{A}

\onepicture[w=5.6cm]{figs/21.png}

\fourchoices[answer=A]
{$Q_{1}=Q_{2}=Q_{3}=Q_{4}$}
{$Q_{1}=Q_{2}=2Q_{3}=2Q_{4}$}
{$2Q_{1}=2Q_{2}=Q_{3}=Q_{4}$}
{$Q_{1}\neq Q_{2}=Q_{3}\neq Q_{4}$}

%% answer: A
%% memo:
\memoanswer{
因 $Q=It=\frac{Bl_{1}v}{R}\cdot\frac{l_{2}}{v}=\frac{BS}{R}$，因此通过电阻 $R$ 的电量与导体棒移动的速度无关，只与磁通量的变化量有关，在四个过程中，回路中磁通量的变化量相等，所以在上述过程中通过电阻 $R$ 的电量都相等，故 A 正确。
}

\item
%% number: 22
%% paperName: 2006年全国理综Ⅰ第 22 题 12 分
%% typeId: 4
%% type: 实验题
%% chapter: 电路
%% point: 电路实验
%% method: 无
%% score: 12
%% degree: 450
%% duplicateId: 0
%% body:
\begin{enumerate}
	\item
	利用图中装置研究双缝干涉现象时，有下面几种说法：

	\onepicture[w=9.0cm]{figs/22.png}

	（A）将屏移近双缝，干涉条纹间距变窄

	（B）将滤光片由蓝色的换成红色的，干涉条纹间距变宽

	（C）将单缝向双缝移动一小段距离后，干涉条纹间距变宽

	（D）换一个两缝之间距离较大的双缝，干涉条纹间距变窄

	（E）去掉滤光片后，干涉现象消失

	其中正确的是 \tkanswer[2.5cm]{ABD}。

	\item
	现要测量某一电压表 V 的内阻。给定的器材有：待测电压表 V（量程 $2\UV$，内阻约为 $4\UkO$）；电流表 mA（量程 $1.2\UmA$，内阻约 $500\UO$）；直流电源 $E$（电动势约 $2.4\UV$，内阻不计）；固定电阻 3 个：$R_{1}=4000\UO$，$R_{2}=10000\UO$，$R_{3}=15000\UO$；电键 S 及导线若干。

	要求测量时两电表指针偏转均超过其量程的一半。

	i. 试从 3 个固定电阻中选用 1 个，与其它器材一起组成测量电路，并在虚线框内画出测量电路的原理图。（要求电路中各器材用题中给定的符号标出。）

	ii. 电路接通后，若电压表读数为 $U$，电流表读数为 $I$，则电压表内阻 $R_{V}=$ \tkanswer[3cm]{}。
\end{enumerate}

%% answer:
\jdanswer{
\begin{enumerate}
	\item
	ABD。
	\item
	i. 电路如图。

	\onepicture[w=6.0cm]{figs/22_an.png}

	ii. $R_{V}=\frac{R_{1}U}{R_{1}I-U}$
\end{enumerate}
}

%% memo:
\memoanswer{
（1）由 $\Delta x=\frac{l}{d}\lambda$ 可知，将屏移近双缝时，$l$ 变小，干涉条纹间距变窄，故 A 正确；蓝光的波长小于红光的波长，将滤光片换成红色后，干涉条纹间距变宽，故 B 正确；改变单缝与双缝间的距离，干涉条纹间距不变，故 C 错误；将双缝间的距离增大时，干涉条纹间距变窄，故 D 正确；去掉滤光片后，干涉现象仍存在，为彩色条纹，故 E 错误。

（2）此题要注意电压表和电流表的量程和测量时两电表指针偏转均要超过其量程的一半这一条件。电压表的额定电流约 $I_{V}=\frac{U}{R_{V}}=\frac{2}{4000}\UA=0.5\times10^{-3}\UA=0.5\UmA$，所以要给电压表并联一个固定电阻，该固定电阻的电阻值约为 $R_{x}=\frac{U}{I_{A}-I_{V}}=\frac{2}{(1.2-0.5)\times10^{-3}}\UO\approx3\times10^{3}\UO$，故固定电阻应选 $R_{1}$，实验电路如图所示。流过定值电阻的电流大小为 $\frac{U}{R_{1}}$，则流过电压表的电流为 $I-\frac{U}{R_{1}}$，故电压表的内阻为 $R_{V}=\frac{U}{I-\frac{U}{R_{1}}}=\frac{UR_{1}}{IR_{1}-U}$。
}

\item
%% number: 23
%% paperName: 2006年全国理综Ⅰ第 23 题 16 分
%% typeId: 6
%% type: 计算题
%% chapter: 直线运动
%% point: 匀速直线运动
%% method: 近似与估算
%% score: 16
%% degree: 450
%% duplicateId: 0
%% body:
天空有近似等高的浓云层。为了测量云层的高度，在水平地面上与观测者的距离为 $d=3.0\Ukm$ 处进行一次爆炸，观测者听到由空气直接传来的爆炸声和由云层反射来的爆炸声时间上相差 $\Delta t=6.0\Us$。试估算云层下表面的高度。已知空气中的声速 $v=\frac{1}{3}\Ukms$。

%% answer:
\jdanswer{$h=2.0\times10^{3}\Um$}

%% memo:
\memoanswer{
如图，A 表示爆炸处，O 表示观测者所在处，$h$ 表示云层下表面的高度。用 $t_{1}$ 表示爆炸声直接传到 O 处所经时间，则有

$d=vt_{1}$ ①

用 $t_{2}$ 表示爆炸声经云层反射到达 O 处所经历时间，因为入射角等于反射角，故有

$2\sqrt{\left(\frac{d}{2}\right)^{2}+h^{2}}=vt_{2}$ ②

已知 $t_{2}-t_{1}=\Delta t$ ③

联立①②③式，可得 $h=\frac{1}{2}\sqrt{(v\Delta t)^{2}+2dv\Delta t}$ ④

代入数值得 $h=2.0\times10^{3}\Um$。

\onepicture[w=5.0cm]{figs/23_an.png}
}

\item
%% number: 24
%% paperName: 2006年全国理综Ⅰ第 24 题 19 分
%% typeId: 6
%% type: 计算题
%% chapter: 运动和力
%% point: 牛顿定律的应用
%% method: 无
%% score: 19
%% degree: 450
%% duplicateId: 0
%% body:
一水平的浅色长传送带上放置一煤块（可视为质点），煤块与传送带之间的动摩擦因数为 $\mu$。初始时，传送带与煤块都是静止的。现让传送带以恒定的加速度 $a_{0}$ 开始运动，当其速度达到 $v_{0}$ 后，便以此速度做匀速运动。经过一段时间，煤块在传送带上留下了一段黑色痕迹后，煤块相对于传送带不再滑动。求此黑色痕迹的长度。

%% answer:
\jdanswer{$\frac{v_{0}^{2}(a_{0}-\mu g)}{2\mu a_{0}g}$}

%% memo:
\memoanswer{
根据“传送带上有黑色痕迹”可知，煤块与传送带之间发生了相对滑动，煤块的加速度 $a$ 小于传送带的加速度 $a_{0}$。根据牛顿定律，可得

$a=\mu g$ ①

设经历时间 $t$，传送带由静止开始加速到速度等于 $v_{0}$，煤块则由静止加速到 $v$，有

$v_{0}=a_{0}t$ ②，$v=at$ ③

由于 $a<a_{0}$，故 $v<v_{0}$，煤块继续受到滑动摩擦力的作用。再经过时间 $t'$，煤块的速度由 $v$ 增加到 $v_{0}$，有

$v_{0}=v+at'$ ④

此后，煤块与传送带运动速度相同，相对于传送带不再滑动，不再产生新的痕迹。

设在煤块的速度从 0 增加到 $v_{0}$ 的整个过程中，传送带和煤块移动的距离分别为 $s_{0}$ 和 $s$，有

$s_{0}=\frac{1}{2}a_{0}t^{2}+v_{0}t'$ ⑤

$s=\frac{v_{0}^{2}}{2a}$ ⑥

传送带上留下的黑色痕迹的长度

$l=s_{0}-s$ ⑦

由以上各式得 $l=\frac{v_{0}^{2}(a_{0}-\mu g)}{2\mu a_{0}g}$ ⑧
}

\item
%% number: 25
%% paperName: 2006年全国理综Ⅰ第 25 题 20 分
%% typeId: 6
%% type: 计算题
%% chapter: 电路
%% point: 电容
%% method: 无
%% score: 20
%% degree: 250
%% duplicateId: 0
%% body:
有个演示实验，在上下面都是金属板的玻璃盒内，放了许多锡箔纸揉成的小球，当上下板间加上电压后，小球就上下不停地跳动。现取以下简化模型进行定量研究。

如图所示，电容量为 $C$ 的平行板电容器的极板 A 和 B 水平放置，相距为 $d$，与电动势为 $E$、内阻可不计的电源相连。设两板之间只有一个质量为 $m$ 的导电小球，小球可视为质点。已知：若小球与极板发生碰撞，则碰撞后小球的速度立即变为零，带电状态也立即改变，改变后，小球所带电荷符号与该极板相同，电量为极板电量的 $\alpha$ 倍（$\alpha\ll1$）。不计带电小球对极板间匀强电场的影响。重力加速度为 $g$。

\begin{enumerate}
	\item
	欲使小球能够不断地在两板间上下往返运动，电动势 $E$ 至少应大于多少？
	\item
	设上述条件已满足，在较长的时间间隔 $T$ 内小球做了很多次往返运动。求在 $T$ 时间内小球往返运动的次数以及通过电源的总电量。
\end{enumerate}

\onepicture[w=5.0cm, align=right]{figs/25.png}

%% answer:
\jdanswer{
\begin{enumerate}
	\item
	$E>\sqrt{\frac{mgd}{\alpha C}}$
	\item
	$n=\frac{T}{\sqrt{\frac{2md^{2}}{\alpha CE^{2}+mgd}}+\sqrt{\frac{2md^{2}}{\alpha CE^{2}-mgd}}}$，$Q'=\frac{2\alpha CET}{\sqrt{\frac{2md^{2}}{\alpha CE^{2}+mgd}}+\sqrt{\frac{2md^{2}}{\alpha CE^{2}-mgd}}}$
\end{enumerate}
}

%% memo:
\memoanswer{
（1）用 $Q$ 表示极板电荷量的大小，$q$ 表示碰后小球电荷量的大小。要使小球能不停地往返运动，小球所受的向上的电场力至少应大于重力，则 $q\frac{E}{d}>mg$ ①，其中 $q=\alpha Q$ ②，又有 $Q=CE$ ③，由①②③式得 $E>\sqrt{\frac{mgd}{\alpha C}}$ ④。

（2）当小球带正电时，小球所受电场力与重力方向相同，向下做加速运动。以 $a_{1}$ 表示其加速度，$t_{1}$ 表示从 A 板到 B 板所用的时间，则有 $q\frac{E}{d}+mg=ma_{1}$ ⑤，$d=\frac{1}{2}a_{1}t_{1}^{2}$ ⑥。当小球带负电时，小球所受电场力与重力方向相反，向上做加速运动，以 $a_{2}$ 表示其加速度，$t_{2}$ 表示从 B 板到 A 板所用的时间，则有 $q\frac{E}{d}-mg=ma_{2}$ ⑦，$d=\frac{1}{2}a_{2}t_{2}^{2}$ ⑧。小球往返一次共用时间为 $t_{1}+t_{2}$，故小球在 $T$ 时间内往返的次数 $n=\frac{T}{t_{1}+t_{2}}$ ⑨。由以上关系式得 $n=\frac{T}{\sqrt{\frac{2md^{2}}{\alpha CE^{2}+mgd}}+\sqrt{\frac{2md^{2}}{\alpha CE^{2}-mgd}}}$ ⑩。小球往返一次通过的电量为 $2q$，在 $T$ 时间内通过电源的总电量为 $Q'=2qn$，由上式及 $n$ 的表达式可得 $Q'=\frac{2\alpha CET}{\sqrt{\frac{2md^{2}}{\alpha CE^{2}+mgd}}+\sqrt{\frac{2md^{2}}{\alpha CE^{2}-mgd}}}$。
}

\end{enumerate}

\end{document}
